\section{Discussion}\label{sec:discussion}

\subsection{Detection Performance}\label{sec:discussion-detection}

On MIO-TCD, where human labels are available, fine-tuning yields a substantial
improvement on the evaluated dataset: mAP@0.5 rises from 0.430 to 0.720 and the
gain is stable across scene groups. On the LTA images the reference is the
YOLO26x teacher, and the joint model reaches 0.638 at the batch-inference
resolution, close to the pretrained model's 0.660 and well above MIO-only
training at 0.201. Joint training improves agreement on the Singapore validation
set while retaining the accuracy gained from MIO-TCD. The human-labeled MIO-TCD
test set provides the accuracy benchmark, and the LTA figures measure
cross-domain agreement with the teacher.

\subsection{Camera Filter and Congestion Levels}\label{sec:discussion-congestion}

The camera filter shows that fixed scenes allow a simple, transparent correction
of a learned detector: it removes 2.0\% of boxes and raises precision by 1.7
points on unseen human-labeled images, without retraining. The gain is small
because most errors on LTA frames are missed vehicles, which a filter cannot
recover; recall, not precision, limits counting on these cameras.

For congestion, a stacked combination of a trained detector, classical image
processing and a zero-shot vision-language model agrees with 74.3\% of held-out
human labels, and recognizes the congested level with precision 0.88 and recall
0.84. Its 2.0-point gain over the detector count lies within the uncertainty of a
400-frame test, so the result supports the method as a working estimate rather
than showing that combining voters is necessary. The main difficulty is the
moderate level, which is usually rated free-flowing, the boundary where the
visual difference between levels is smallest. Because every design choice was
made on practice labels and the test labels were used once, 74.3\% is a fair
estimate for new frames from the same cameras and period. The levels are a
visual judgement of the outlined road, not a calibrated flow or speed, so they do
not close the integration gap described in Section~\ref{sec:method-integration}.

\subsection{Forecasting Performance}\label{sec:discussion-forecasting}

The forecasting module was selected through the common-protocol comparison.
GWNET is the best-performing model under the evaluated setting. Across five
paired seeds, no clear difference was established between MultiHeadSTGCN and
MTGNN or AGCRN, and the limited number of seeds does not support an equivalence
claim. The backbone-matched configuration lowers MAE relative to STGCN by about
14.8\% in the staged ablation and by 14.9\% in the default-configuration
benchmark, but these comparisons also change inputs, loss and training strategy;
they should not be read as an isolated architectural effect. Under both native
and unified inputs, GWNET retains the lowest MAE in this benchmark.

\subsection{Module Integration}\label{sec:discussion-integration}

The detection and forecasting modules are part of one prototype and operate on
different data domains. The detection module consumes camera images; the
forecasting module consumes road-network time series. The workbench reads both
sets of outputs and labels their sources, so a user can inspect the monitoring
view and the forecast view with the corresponding data source identified. The
separation follows the available data coverage; both modules remain part of the
same prototype. Section~\ref{sec:method-integration} defines the interface for a
connected system, and Section~\ref{sec:discussion-roadmap} discusses the
implications of moving from the prototype to that system.

\subsection{Implications for Future Integration}\label{sec:discussion-roadmap}

Section~\ref{sec:method-integration} defines the data interface; the main
implication is that integration is a measurement problem before it is a modeling
problem. Camera counts must be calibrated against a trusted volume source and
aggregated to a fixed interval, because a detector count per still image is not a
flow series. The same locations also need time-aligned flow, and ideally
occupancy and speed, over weeks or months. Where the existing view is too distant
or poorly lit, close-range cameras or dedicated sensors are the practical way to
obtain those measurements.

A second implication concerns model transfer. A forecasting model trained on
PEMSD8 learns California sensor dynamics; applying it to a Singapore series
would require retraining and a new evaluation on the local split. Detection
counts could enter as an exogenous variable, but only after calibration and
timestamp alignment. These steps define the work that remains before the
workbench can present one local monitoring-and-forecasting workflow rather than
two source-labeled views.

\subsection{Limitations and Scope}\label{sec:discussion-limitations}

The reported system is a time-bounded prototype and pipeline design. It
demonstrates the processing chain and evaluation protocol. An operational
LTA-to-Singapore system would require calibrated sensing and longer data
collection.

\begin{itemize}
  \item LTA results use the YOLO26x pseudo-labels as reference, so they measure
        agreement with the teacher, and the student inherits the teacher's
        missed motorcycles and false positives. Human-labeled MIO-TCD provides
        the accuracy benchmark.
  \item Nighttime LTA frames are difficult to annotate even for humans:
        vehicles are small, poorly lit and often blurred. Teacher labels and
        human labels collected at night are the least reliable reference.
  \item LTA agreement depends on inference resolution: the joint model leads at
        960 pixels and the pretrained model at 1280 pixels.
  \item Scene groups provide a content-based proxy for camera identity, so the
        MIO-TCD test set can retain residual same-camera overlap. In an audit of
        40 test images, each paired with its most similar training image, a
        preliminary visual review found 9 same-camera pairs. The test score
        therefore describes a mix of seen and unseen cameras, not generalization
        to new cameras.
  \item MIO-TCD is evaluated on a filtered four-class subset with a
        project-defined class mapping.
  \item The workbench risk ranking is a threshold-based indicator defined on
        predicted speed. It supports comparative inspection but has not been
        validated against an independent congestion label.
  \item The camera filter and congestion evaluations rely on a single
        annotator. The congestion test set has 50 frames per camera, and two
        weeks give each weekday only two or three days, so per-camera and
        per-slot figures are uncertain.
  \item Congestion is judged only inside the calibrated detection zones, not
        along the whole road, and the zones and road regions must be
        recalibrated if a camera is moved or re-aimed.
  \item The LTA images are still images at ten-minute intervals. They support
        vehicle counts per frame but not continuous passage counts at a point.
        A PEMSD8-style flow series would require continuous video, tracking,
        lane or region calibration, and validation against a loop or radar
        sensor.
  \item Full local integration requires additional IoT measurements (loop
        detectors, radar, LiDAR or traffic counters) at the camera locations to
        provide calibrated flow, occupancy and speed. Close-range road cameras
        are an alternative when the existing camera geometry or lighting is
        unfavorable.
\end{itemize}
